\Cref{sec:results} found the measured degrees obeying the laws of the formulation chapters within their stated scope, and departing from them on singular problems (\Cref{def:regular}), on two axes and on a lift past its boundary. Exponents record how fast a degree grows; they do not explain why carrying the derivatives lowers the Hamiltonian's norm and endangers its gap by one and the same act, nor what a lower degree is worth once gates, amplification, readout and the classical alternatives are counted, nor how far a study made entirely of classical computation can vouch for a quantum algorithm. This chapter asks \emph{what mechanism produces the descriptor's advantage and its limits, and what the measurements imply beyond the benchmark problems.} The answer it develops is that ground-state preparation of a differential equation costs what the conditioning of its residual costs, and that the formulation, rather than the circuit, decides that conditioning.

\subsection{Conditioning Moves from the Norm to the Gap}
\label{sec:disc-tradeoff}

The descriptor's advantage belongs to the operator and not to its encoding. The structured encoding holds the standard form's subnormalisation within a constant of its norm, and the degree still grows as $N^{2k}$ (\Cref{tab:growth}), so the exponent must be read from the residual itself. The factorisation $\G=\Bt^{-1}\Dt$ of \Cref{eq:G-banded} shows where it comes from. Of the two powers of $N$ in $\|\G\|_2$, one is $\|\Dt\|_2=2(N-1)$ and the other belongs to the inverse, $\Bt^{-1}$ being a triangle of ones on each parity class. Eliminating $k$ derivatives composes $\G$ $k$ times, so $\|R_{\mathrm{std}}\|_2=\Theta(N^{2k})$, and the Gram product doubles the exponent again. The descriptor keeps the banded factors, which avoids the inverse, and carries each derivative as its own unknown, which avoids the composition: every recurrence row applies $\Dt$ and $\Bt$ once and uninverted, no block exceeds $\mathcal O(N)$ in norm at fixed $\zeta$ whatever the order, and $\|\Hdesc\|=\Theta(N^2)$ (\Cref{sec:descriptor-spectrum}).

What is removed from the norm reappears in the gap. By \Cref{eq:descriptor-quotient} the descriptor measures a state's residual against $\sum_m\|\G^mp\|^2$, the very quantities the standard form pays for in its norm, so a state whose carried derivatives are large against its residual has a small Rayleigh quotient. The reformulation therefore relocates the conditioning of differentiation rather than removing it: from a norm that grows with $N$ on every problem, to a gap that falls only where the equation fails to control one of the derivatives it carries. \Cref{thm:flat-gap} proves that the control holds on a regular problem, and both ways in which the gap was measured to fall are failures of that control. Where the leading coefficient vanishes, a state concentrated near the zero has a top derivative that the equation does not penalise; on two axes the equation constrains only a sum of top derivatives, so a state whose second derivatives along the two axes nearly cancel meets it with a small residual (\Cref{sec:descriptor-gap}). The stiff R2c shows how severe the erosion can be even on a regular problem, its unscaled gap at $n=8$ being $\res[2,sci]{build/R2c/descriptor/n8/unscaled}{gap}$ against $\res[3]{build/R2c/descriptor/n8}{gap}$ rescaled (\Cref{sec:results-gap}). The mass form of a first-order equation removes the inverse without carrying any derivative, so its states are measured by the field alone, which is consistent with its gap exceeding a chain's on R7 (\Cref{sec:systems-fields}). At higher order no such option exists: the chain lowers the norm only by carrying the derivatives, and so only at a risk to the gap.

Classical numerical analysis met the same trade from two directions. First-order-system least squares carries the derivatives so as not to square a high-order operator's condition number, and its theory rests on norm equivalence, the requirement that the least-squares functional bound every carried unknown in its own norm \cite{cai1994fosls}. \Cref{thm:flat-gap} is a discrete, one-axis instance of the lower half of that equivalence, coercivity, and the collapse on two axes is what its failure would look like. The datum rows of \Cref{sec:standard-constraints} weigh a boundary function equally at every frequency, which may be too weak to control the second derivatives of a harmonic or diffusive state of high frequency. That reading is inferred from the mechanism and not measured, and it makes datum rows weighted towards high frequencies the first remedy to test (\Cref{sec:future-work}). The ultraspherical method reaches the same banded factors and finds its systems well conditioned once diagonally preconditioned \cite{olver2013}, which is the role block rescaling takes here. The increment of this thesis is therefore not the idea of carrying derivatives through banded factors but its transfer to a block-encoded Hamiltonian, where conditioning is paid in queries, together with the measurement of what it costs and where it fails.

The price that remains consists of factors independent of $N$, and they decide whether the reformulation pays at small resolution. The block register multiplies the state by $2^{n_f}$, four blocks at second order and eight for Laplace's and the wave equation (\Cref{sec:app-variable-coeff}). A query applies more terms through that register, so on constant-coefficient problems it costs more two-qubit gates than the structured standard form's at the sizes built and verified, $\res{encode/R2a/descriptor/n2}{cx_count}$ against $\res{encode/R2a/standard/n2}{cx_count}$ on R2a at $n=2$, the two converging to $\res{query/stateprep/R2a/descriptor/n8/T3/filter=minimax}{gates}$ against $\res{query/stateprep/R2a/standard/n8/T3/filter=minimax}{gates}$ at $n=8$ (\Cref{sec:results-scaling}). Block rescaling adds no ancilla, but it lengthens the Pauli expansion of the register's patterns by a factor fixed by the equation. Whether the reformulation pays at all is decided instead by the exponents, and they depend on the class of problem: on a singular problem the rule's $\zeta\propto N$ raises the descriptor's norm to $\Theta(N^4)$ at second order, and on two axes or at a double zero no block weight restores the gap (\Cref{sec:descriptor-rescaling}).

\subsection{End-to-End Complexity}
\label{sec:disc-complexity}

Neither form is polylogarithmic in $N$. The degree of \Cref{eq:degree}, $d=\Theta\bigl(\sqrt{\alpha_H/\Delta}\,\log(1/\varepsilon)\bigr)$ with $\alpha_H\ge\|H\|$ for any block encoding, is the residual's effective condition number times $\log(1/\varepsilon)$. It has the form of the optimal linear-systems law, $\mathcal O(\kappa\log(1/\varepsilon))$ queries \cite{costa2022optimal}, which no algorithm improves to $\kappa^{1-\delta}$ unless $\mathrm{BQP}=\mathrm{PSPACE}$ \cite{harrow2009hhl}, so the ground-state route does not escape conditioning; it meets it in another guise. Read this way, the standard form's $\tilde\Theta(N^{2k})$ for constant coefficients with the best exact encoding (\Cref{eq:standard-wall}) is the conditioning of an order-$k$ least-squares discretisation, and the descriptor's $\tilde\Theta(N)$ on a regular problem (\Cref{eq:descriptor-degree}) is that of a first-order one, the conditioning that the linear-systems solvers for ordinary equations also pay by reducing the order before they discretise (\Cref{sec:bg-de-linear}). On singular problems and on two axes the descriptor's degree is about $\tilde\Theta(N^2)$ or steeper. Without a source, or with a polynomial source of fixed degree, a query costs $\mathcal O(n^2)$ two-qubit gates in both forms (\Cref{prop:standard-cost,prop:descriptor-cost}), a general source adding $\mathcal O(N)$, and amplification multiplies the queries by $2k_{\mathrm{aa}}+1=\mathcal O(1/\gamma)$ (\Cref{eq:total-cost}). On a regular problem with constant coefficients a preparation therefore costs $\tilde\Theta(N\,\mathrm{poly}(n)/\gamma)$ gates in the descriptor against $\tilde\Theta(N^{2k}\,\mathrm{poly}(n)/\gamma)$ in the standard form. The overlap $\gamma$ is measured rather than bounded, and it differs between the forms, the descriptor's being capped by the field block's weight, which the rescaling restores (\Cref{sec:results-gap}).

The measured degrees show how the two costs hidden in $\alpha_H$ divide. On R2a at $n=8$ Wu et al.'s encoding implies a degree of $\res[2,sci]{stateprep/R2a/standard/n8/T3/filter=minimax/encoder=published}{degree}$, and an exact structured encoding of the same residual one of $\res[2,sci]{stateprep/R2a/standard/n8/T3/filter=minimax}{degree}$: the larger factor in the degree their circuits imply is the encoder's, and it is removable. The remainder, $\res[2,sci]{stateprep/R2a/standard/n8/T3/filter=minimax}{ideal_degree}$ at the ideal bound, is the Hamiltonian's, against the descriptor's $\res{stateprep/R2a/descriptor/n8/T3/filter=minimax}{degree}$ as built and $\res{stateprep/R2a/descriptor/n8/T3/filter=minimax}{ideal_degree}$ at the ideal bound (T3, \Cref{tab:headline}). Because the descriptor's encoding is the tighter, within $\res[3]{looseness/stateprep/R2a/descriptor/n8/T3/filter=minimax}{alpha_ratio}$ of its residual's norm against $\res[3]{looseness/stateprep/R2a/standard/n8/T3/filter=minimax}{alpha_ratio}$ for the structured encoding, the as-built comparison favours the descriptor, and only the ideal bound compares the two Hamiltonians. For readers of Wu et al.\ the consequence is that the method's cost lives in the Hamiltonian it constructs and not in the circuits it reports. A gate count per query, which is how they price differentiation, is accurate and incomplete, because the cost of differentiation returns in the number of queries; the separation of the encoder's cost from the Hamiltonian's is an audit that any filter-based algorithm on a block-encoded operator admits.

How far the floor extends is a separate question. \Cref{eq:standard-wall} holds for every block encoding of $R_{\mathrm{std}}$ under any polynomial filter of the walk, since by Chebyshev's extremal property none suppresses the excited spectrum at a lower degree than the minimax filter (\Cref{sec:bg-filtering}). Over all algorithms the evidence is weaker. Lin and Tong prove, by reduction from unstructured search, that their ground-state algorithm's dependence on the gap and the overlap is essentially optimal for an oracle that block-encodes $H$ \cite{lintong2020groundstate}. The walk here is built from an encoding of $R$, which supplies the square root in \Cref{eq:degree}, so their bound does not apply as it stands, and whether their hard instances arise as stacked residuals is not shown. The thesis therefore claims the floor for polynomial filters of the walk, and for Wu et al.'s second-order examples with constant coefficients it follows that $\|\Hstd\|=\Theta(N^8)$ sets the number of queries under any block encoding, whatever circuit implements one query (\Cref{sec:standard-hamiltonian}). The descriptor lowers the floor by replacing $R$, which no statement about the encodings of a fixed $R$ excludes.

Within that family the choice of filter sets the constant and leaves the exponent alone. The \gls{gqsp} circuit realises the minimax edge filter and Wu et al.'s imaginary-time filter from the same walk at the same normalisation (\Cref{sec:bg-qsp,app:gqsp-phases}), so their comparison isolates the polynomial. At equal suppression imaginary time needs $\res[3]{filters/synthetic/edge}{low}$ to $\res[3]{filters/synthetic/edge}{high}$ times the minimax degree on synthetic spectra, whatever the gap, and the ratio does not drift with the resolution, $\res{filters/R1/descriptor/n2}{e6_qite_worst}$ against $\res{filters/R1/descriptor/n2}{e6_minimax_certified}$ on R1 at $n=2$ and $\res{filters/R1/descriptor/n8}{e6_qite_worst}$ against $\res{filters/R1/descriptor/n8}{e6_minimax_certified}$ at $n=8$ (T2, \Cref{app:filter-study}). The reverse cases are instructive. At equal measured infidelity and a loose target, $\varepsilon=10^{-4}$, imaginary time needs fewer degrees on three of the eight cases, $\res{filters/R3a/descriptor/n3}{e4_qite_degree}$ against $\res{filters/R3a/descriptor/n3}{e4_minimax_degree}$ on R3a in the descriptor, because its exponential suppresses each excited state more the higher it lies, while the minimax filter holds every one at the equiripple level. That dividend comes from the structure of the spectrum above the edge, which an a priori pipeline neither knows nor can certify; the minimax filter is optimal for exactly the worst case that the edge alone determines (\Cref{prop:minimax-optimal}), and so it is the filter that answer-independent parameters can guarantee. The lower degree also has a numerical benefit: three of the reproduction's imaginary-time circuit simulations, at degrees of order $10^5$, miss the agreement bar through rounding accumulated over the walks, while every minimax record meets it (\Cref{app:repro-t1}). The filter is thus an engineering constant worth a factor below two, and Wu et al.'s open question about the gap was the right one to ask: the exponent lives in the residual.

What $N$ costs depends on what fixes it, and here the Chebyshev basis is decisive. Solved classically through its reference, the standard form converges spectrally, its error falling to $\res[2,sci]{build/R2a/standard/best}{field_err_linf}$ on R2a at $n=\res{build/R2a/standard/best}{n}$ and to $\res[2,sci]{build/R5b/standard/best}{field_err_linf}$ on the heat equation at $n=\res{build/R5b/standard/best}{n}$ per axis, as the geometric decay of an analytic solution's coefficients predicts \cite{trefethen2013approximation}. Beyond that resolution its error rises again, to $\res[2,sci]{build/R2a/standard/n10}{field_err_linf}$ at $n=10$, where the descriptor's, with the same kernel, is $\res[2,sci]{build/R2a/descriptor/n10}{field_err_linf}$ (\Cref{sec:results-faithful}). The rise is consistent with rounding amplified by a residual whose norm grows as $N^4$, so the conditioning that sets the quantum degree is already visible in the classical solve, a diagnostic that needs no quantum computer. The prepared state shows the same saturation from the other side: on R2a in the descriptor the decoded field's error is $\res[2,sci]{stateprep/R2a/descriptor/n4/T2/filter=minimax}{field_err_linf}$ at $n=4$ and $\res[2,sci]{stateprep/R2a/descriptor/n8/T2/filter=minimax}{field_err_linf}$ at $n=8$, set by $\varepsilon$ while the discretisation's own error falls to $\res[2,sci]{stateprep/R2a/descriptor/n8/T2/filter=minimax}{eps_disc}$, so resolution beyond what the solution needs buys only degree. For an analytic solution a field accurate to $\delta$ needs $N$ growing as $\log(1/\delta)$, both degrees are then polylogarithmic in the precision, and the comparison of exponents in $N$ implies no speed-up. On one axis a degree linear in $N$ is no asymptotic advantage in any case: the ultraspherical method solves such an equation classically in $\mathcal O(N)$ operations once its coefficients are resolved \cite{olver2013}, and decoding the whole field from a quantum state needs of order $N$ estimates (\Cref{sec:bg-filtering}). An advantage could appear only on many axes, where a classical solver's cost grows at least as the number of unknowns, $N^D$ on $D$ axes. That is precisely where the descriptor's gap falls: on the two axes measured its estimated degree grows as $N^{\res[2]{stateprep/R5b/descriptor/growth/T3/filter=minimax}{degree}}$ per axis on the heat equation and $N^{\res[2]{stateprep/R5a/descriptor/growth/T3/filter=minimax}{degree}}$ on Laplace's (\Cref{tab:growth}), faster than the $N^2$ unknowns, though between the least settled resolutions of the study. Whether a residual can keep its gap on several axes therefore decides whether this route can compete with classical spectral solvers at all, and the thesis leaves that question open.

\subsection{Where the Construction Applies}
\label{sec:disc-applicability}

\Cref{tab:applicability} turns the measured classes into guidance a practitioner can act on. Its organising principle follows from \Cref{sec:disc-tradeoff}: the descriptor pays where the equation controls every derivative it carries, and on nonlinear problems only while the lift converges.

\begin{table}[htbp]
    \centering
    \caption[Where the descriptor form pays, by class of problem]{\textbf{The descriptor pays where the equation controls every derivative it carries, and on nonlinear problems only while the lift converges.} For each class of problem, the benchmark problems that represent it, the descriptor's gain over the standard form in the filter's degree, and the choice our evidence supports. $k$, order; regular, singular: one axis, the leading coefficient without and with a zero on $[-1,1]$ (\Cref{def:regular}); $n_t$, time-axis qubits; ideal bound, $\alpha_R=\|R\|_2$. Gains are proved (\Cref{thm:flat-gap}) or measured as resource estimates (T3, \Cref{tab:growth,tab:catalogue-growth}) between stated resolutions; the complex-valued row is a statement about the formulation, with no complex residual prepared.}
    \label{tab:applicability}
    \footnotesize
    \renewcommand{\arraystretch}{1.15}
    \begin{tabularx}{\linewidth}{@{}>{\RaggedRight}p{0.2\linewidth} >{\RaggedRight}X >{\RaggedRight}X@{}}
        \toprule
        Class (problems) & The descriptor's gain & What to do \\
        \midrule
        Regular (R1, R2, R4b--d) & degree $\tilde\Theta(N)$ against $\tilde\Theta(N^{2k})$ for constant coefficients; gap proved $\Theta(1)$ & the chain with $\zeta$ from \Cref{eq:zeta-rule} for $k\ge2$, the mass form for $k=1$; expect more gates per query at small $n$ \\
        Singular, simple zero (R3, R4a) & degree about $\tilde\Theta(N^2)$; against $N^4$ where the zero is at one end (R4a), none at the ideal bound where it is at both (R3) & $\zeta=N/2$; judge at the ideal bound, since as built the gain on R3 is the encoder's \\
        Singular, double zero & none claimed: balancing needs a norm above the standard form's own & substitute $f=\sqrt{1-x^2}\,v$ to leave simple zeros, as for R3c and R3d \\
        Two axes, no mixed derivative (R5) & degree's exponent per axis below the standard form's; reversed on the heat equation at $n=4$ per axis; gap falls even when rescaled & $\zeta_a=N_a$; budget a state of four to eight blocks; compare at the ideal bound and expect no guaranteed gain \\
        First-order systems (R7) & regular: about $N$ against $N^2$ & the mass form, with no block register and a larger gap than a chain's; data as ratio rows; the coupling as matchings \\
        Complex-valued linear equations & that of the real system after realification (R7 is $z'=-\mathrm{i}\pi z$ realified) & realify, at one more field-register qubit; native complex preparation is not built \\
        Polynomial nonlinearity, lift converging (R8; R9 to $a/\nu=3$; R10 to $a/\nu=10$, and at $30$ on $12$ modes) & through the time axis only, from $n_t=4$, narrowing as the modes grow & the Carleman lift, mass form in time; raise $n_t$ with the truncation; measure the lift's error per order; never the doubled space, which projects \\
        Past the lift's boundary (R9 at $a/\nu=10$; R10 at $a/\nu=30$ from $28$ modes) & none: the kernel stays trusted while the lift diverges & shorten the horizon or lower the strength; read convergence from the lift's error, never from trust \\
        \bottomrule
    \end{tabularx}
\end{table}

The split in the simple-zero row turns on where the coefficient vanishes. The standard form's norm is set by the top Chebyshev modes, whose derivatives are largest at the ends of the interval by Markov's inequality \cite{borweinerdelyi1995}. A coefficient that vanishes at both ends, $1-x^2$ on R3, suppresses those modes where they peak, and there $\|\Hstd\|$ grows only as $N^4$; $x-1$ on R4a vanishes at one end, and the $N^8$ law holds (\Cref{sec:results-scaling}). The rescaled descriptor's norm grows as $N^4$ on both, so it gains on R4a and not on R3. The account is an inference from two problems, consistent with \Cref{sec:standard-hamiltonian} but not a theorem. It makes a testable prediction that no benchmark problem checks: an interior zero, which suppresses neither end, should leave the standard form's $N^{4k}$ law, and with it the descriptor's advantage, in place.

On the nonlinear rows the guidance follows $a/\nu$ rather than $\rho$, and the reason is physical. The ratio $\rho=a\|F_2\|/\nu$ is a worst case over all data of norm $a$, and $\|F_2\|$ grows with the modes because the advection term couples into a mode with a strength that rises with its wavenumber (\Cref{sec:fluids}). Burgers' data $a\sin x$ occupy one mode and reach the others only through the nonlinearity, so at moderate strength the high modes, where $F_2$ is largest, plausibly carry too little of the solution to move the lift's error, which would explain why that error is the same on four modes as on sixteen (\Cref{sec:results-nonlinear}). The vortex pair puts several wavenumbers into its data, its triads are active from the start, and there the error does grow with the modes (\Cref{sec:results-ns}). Both readings are inferences from the measured errors. They agree with $\rho<1$ being sufficient and far from necessary, and they explain why a worst-case ratio is a poor guide to a particular flow.

The same rows carry the thesis's sharpest caution. A trusted kernel certifies that the lifted linear problem has an isolated ground state, which is what the filter needs; it says nothing about whether that ground state is the nonlinear solution. On Burgers at $a/\nu=10$ both forms prepare the lift to $\varepsilon$, and the decoded field is wrong by $\res[3]{stateprep/R9e-M4/descriptor/n5/T2/carleman=5/filter=minimax}{field_err_linf}$ against an amplitude of $10$ (T2); on Navier--Stokes at $a/\nu=30$ from $28$ modes the lift diverges with every kernel trusted. In neither case does anything in the residual's spectrum signal the failure (\Cref{sec:results-nonlinear,sec:results-ns}). A spectral certificate is therefore a certificate of the truncation and not of the physics, and every Carleman preparation needs the lift's measured error per order beside it. The shorter window that restores convergence on $28$ modes is what the lift's time-dependent radius of convergence leads one to expect \cite{forets2017carleman}, and it is the practical lever the table recommends.

The time axis decides the descriptor's gain on those rows. At one Carleman order the forms differ in the time derivative, of norm $2(N_t-1)$ in the mass form against $\Theta(N_t^2)$, and share the coupling, which the mass form multiplies by $\|\Bt_\xi\|_2\in[2,2+\sqrt2]$ (\Cref{sec:prep-vs-proj}). On a coarse time axis, at strong forcing and on many modes, the coupling dominates and that factor makes the mass form's norm the larger, which accounts for the standard form's wins at $n_t=3$ on Burgers and for the descriptor's gain from $n_t=4$. The same balance, with the loosening of the descriptor's encoding of the coupling as its matchings multiply, narrows the Navier--Stokes margin as the modes grow at fixed $n_t$ and predicts a reversal there on enough modes, not reached on up to $48$ (\Cref{sec:results-ns}). The account assumes comparable gaps in the two forms and was not checked term by term, but its recommendation, to raise $n_t$ with the truncation, does not depend on that assumption. The reformulation's advantage therefore survives nonlinearity, carried by the one axis on which the lifted system still differentiates.

\subsection{Feasibility and the Reach of Classical Evidence}
\label{sec:disc-feasibility}

Every cost here is that of a fault-tolerant machine (\Cref{sec:bg-qc}). On R2a at $n=8$ the descriptor's preparation is estimated at $\res{stateprep/R2a/descriptor/n8/T3/filter=minimax}{n_qubits}$ logical qubits and $\res[2,sci]{stateprep/R2a/descriptor/n8/T3/filter=minimax}{gate_count}$ two-qubit gates, $\res[1]{stateprep/R2a/descriptor/n8/T3/filter=minimax}{preparations}$ amplified repetitions of a filter of degree $\res{stateprep/R2a/descriptor/n8/T3/filter=minimax}{degree}$, against $\res[2,sci]{stateprep/R2a/standard/n8/T3/filter=minimax}{gate_count}$ for the structured standard form and $\res[2,sci]{stateprep/R2a/standard/n8/T3/filter=minimax/encoder=published}{gate_count}$ with Wu et al.'s encoding (T3, \Cref{tab:headline}). The widths stay modest, the widest circuits estimated, Navier--Stokes on twelve modes at $K=4$ and $n_t=5$, needing $\res{stateprep/R10c-M12/descriptor/n5/T3/carleman=4/filter=minimax}{n_qubits}$ qubits in the descriptor and $\res{stateprep/R10c-M12/standard/n5/T3/carleman=4/filter=minimax}{n_qubits}$ in the standard form; the obstacle is the gate count. The descriptor's count on R2a buys one second-order equation on $N=2^8$ coefficients, a problem of no classical difficulty, and the reformulation divides the count by about the ratio of the two degrees, which grows as $N^{2k-1}$ on such a problem; that alone does not make the problem worth a quantum computer. The estimates count logical two-qubit gates after transpiling to $\{\mathrm{CX},U\}$ (\Cref{sec:pipeline-design}). They omit the non-Clifford gates from which an error-corrected machine synthesises the rotations $U$, whose count dominates its cost, together with the code's overhead and routing, and they record no depth, which grows with the degree because the walks run in sequence. Jennings et al.\ cost a linear-systems route to linear \glspl{ode} in logical qubits and in queries to the block encoding of the system's matrix, with their constant factors \cite{jennings2023cost}; their counts and these meet at the query, whose gates are supplied here for both forms, but their examples are time evolutions of stable systems and are not compared with ours.

Depth rather than width places the method on error-corrected hardware at every size run. A filter of $\res{stateprep/R2a/descriptor/n8/T3/filter=minimax}{degree}$ sequential walks of $\res{query/stateprep/R2a/descriptor/n8/T3/filter=minimax}{gates}$ two-qubit gates each lies far beyond the depth that uncorrected processors execute reliably \cite{preskill2018nisq}, and even the smallest preparation, the decay $f'+2f=0$ (R1) at $n=2$ in the descriptor form, fits on $\res{stateprep/R1/descriptor/n2/T1/filter=minimax}{n_qubits}$ qubits but applies $\res{stateprep/R1/descriptor/n2/T1/filter=minimax}{degree}$ sequential queries of $\res{encode/R1/descriptor/n2}{cx_count}$ two-qubit gates each (circuit simulation). No preparation here is a candidate for a noisy device, and the variational methods that target such devices answer a different question.

The same arithmetic bounds what a classical study can vouch for. The production ran $\res{production}{runs}$ runs on the Setonix supercomputer \cite{pawsey2023setonix}, $\res{production}{hours}$ hours of computation in all, and its ceiling was set by circuit simulation, whose statevector memory doubles with every qubit and whose time grows with the degree (\Cref{app:t1-model}). The widest circuit simulated has $\res{stateprep/R4d/descriptor/n3/T1/filter=minimax}{n_qubits}$ qubits; the narrowest two-axis descriptor circuit, at $\res{stateprep/R5b/descriptor/n2x2/T1/filter=minimax}{n_qubits}$ qubits, estimated its own run at $\res[2]{stateprep/R5b/descriptor/n2x2/T1/filter=minimax}{predicted_seconds}$\,s against a budget of $\res[2]{stateprep/R5b/descriptor/n2x2/T1/filter=minimax}{budget_seconds}$\,s and was refused before it began, one of the $\res{production}{infeasible}$ runs recorded that way (\Cref{sec:results-faithful}). The resolutions at which the method's open question lives, several axes and many modes, thus lie beyond direct simulation on a national facility, a limit of the classical simulator rather than of the circuit. What carries the study past it is the argument that the exact emulation equals the circuit's post-selected output (\Cref{sec:pipeline-verification}), confirmed on all $\res{agreement/T1/minimax}{count}$ minimax circuit simulations to within $10^{-10}$, which licenses the emulation to filter degrees of $2^{20}$ and the resource estimate beyond. For any filter-based algorithm on a block-encoded operator this is the transferable lesson: an emulation proved equal to the circuit and checked against it where both run is far cheaper than a statevector simulation, and it is the only route by which a classical study reaches the resolutions where the scaling matters.

\subsection{Limitations}
\label{sec:limitations}

Every result is a classical computation, and nothing ran on hardware. The claims are cost laws of operators and filters, which classical computation tests exactly; it cannot test noise, and the method's robustness to it is unknown. Exact emulation stands in for the circuit beyond $\res{stateprep/R4d/descriptor/n3/T1/filter=minimax}{n_qubits}$ qubits by an argument that the agreement of the two confirms only where both run, the widest reflections, of the two-axis descriptor problems and the larger lifts, are built and counted but not verified whole, and Navier--Stokes is reached by emulation and estimate alone (\Cref{sec:results-ns}). The growth exponents are two-point slopes between stated resolutions, consistent with the laws rather than fits of them, and the two-axis slopes, between $n=4$ and $5$ per axis, are the least settled. Where $N$ must grow with the problem, as on a stiff equation, the dependence of the gap's constant on the stiffness was measured on R2c alone, and beyond $8{,}192$ columns the reference's error bound is an a posteriori estimate (\Cref{app:iterative-reference}).

Several inputs are supplied by the classical reference rather than found by the method, and this is the limitation with the widest consequence. Every degree uses the spectral edge $\sigma_0,\sigma_1$ of the reference (\Cref{sec:pipeline-filter}). An underestimated gap raises the degree only as the inverse square root of the bound and still meets $\varepsilon$ (\Cref{app:edge-study}), so the laws survive any lower bound of the right order in $N$; \Cref{thm:flat-gap}, proved by compactness, gives that order on regular problems without its constant, which is measured here and supplied nowhere. The exact constraint points of R4 and R6b and the heat equation's initial function at $t=-1$ are set up from the reference solution when the problem is built, an answer-dependent set-up that a new problem must supply for itself (\Cref{sec:pipeline-inputs}), and the two constants in the block weight $\zeta$ were selected by calibration on the benchmark problems and then held fixed, so the weight's check against the best $\zeta$ on a grid is in-sample (\Cref{sec:descriptor-rescaling}). Trust in a kernel and convergence of a lift are likewise judged against references. Where the method would matter, beyond classical reach, none of these references exists, so a quantum preparation would need to certify its own edge and its own convergence. The output, finally, is a state: the field is decoded classically from amplitudes that a quantum computer does not expose, and Wu et al.'s interferometric readout is built and verified, with its correction for post-selection, only at small $n$ (\Cref{sec:pipeline-readout}).

The descriptor's scope ends where its gap does. The gap is proved flat only on regular problems, measured flat after rescaling on singular ones with a simple zero, and falls on two axes; a double zero needs the substitution of \Cref{sec:standard-constraints}, whose form presumes that the zero's order is known from the equation, and mixed derivatives have no block and are not supported (\Cref{sec:descriptor-constraints}). The full ultraspherical reduction, a banded residual on $N$ unknowns with no block register, would remove the inverse without carrying any derivative, but it was neither built nor compared (\Cref{sec:bg-ultraspherical}), and nothing here predicts whether it keeps the descriptor's norm without eroding the gap.

Of the extensions, only first-order systems were run. Complex-valued fields are carried by realification alone; no complex residual was prepared, and the one native complex encoding, the tensor layout's coupling, is verified but prepares nothing (\Cref{sec:systems-complex}). The doubled space projects, so the \gls{dct} that its circuits cite rather than build matters to no preparation (\Cref{sec:pipeline-design}). The thesis has no test of a lift's divergence that reads the residual alone (\Cref{sec:results-nonlinear}). Wu et al.'s partial differential and nonlinear examples, their Figs~6 and~7, were not reproduced under any hypothesis tested, and the authors were not contacted (\Cref{app:repro-unreproduced}); this bears on the audit of the paper rather than on the comparison, since the standard form is specified independently of those panels.

The chapter's findings reduce to one account. The number of queries grows with the resolution as the residual's effective condition number, and in Wu et al.'s circuits the larger part of that growth is the encoder's, which an exact encoding removes, while what remains, $\tilde\Theta(N^{2k})$ for constant coefficients, belongs to the Hamiltonian and survives every block encoding under every polynomial filter of the walk. Reformulating the residual so that it carries its derivatives keeps the kernel and lowers that growth to $\tilde\Theta(N)$ on a regular problem, because it moves the conditioning of differentiation from the norm, where it is paid on every problem, into the gap, where it is paid only when the equation loses control of a carried derivative; block rescaling repairs that loss on one axis without ancilla, and nothing measured here repairs it on two. Nonlinear equations are prepared, not projected, once a Carleman lift replaces the doubled space, and the reformulation's advantage survives there through the time axis, but a prepared lift is certified only as a truncation. The filter, the encoder and the circuit change constants; the formulation changes exponents. Neither form is polylogarithmic, and the two limits that bound every verdict, the gap on several axes and the classical references on which the edge and the lift's convergence rest, are where \Cref{sec:conclusions} takes up the further work.
